% ============================================================
%  IP Policy Communication and Technology Transfer
%  Sharma, Wang, Basnet, Cossco (2026)
%  Working Paper
% ============================================================
\documentclass[12pt]{article}

\usepackage[margin=1.25in]{geometry}
\usepackage{setspace}
\usepackage{amsmath,amssymb}
\usepackage{booktabs}
\usepackage{threeparttable}
\usepackage{tabularx}
\usepackage{multirow}
\usepackage{graphicx}
\usepackage{natbib}
\usepackage[hidelinks]{hyperref}
\usepackage{xcolor}
\usepackage{caption}
\usepackage{subcaption}
\usepackage{footnote}
\usepackage{rotating}

\newcommand{\PCSI}{\textsc{pcsi}}

\setstretch{1.5}
\setlength{\parskip}{4pt}

\newcommand{\sone}{$^{*}$}
\newcommand{\stwo}{$^{**}$}
\newcommand{\sthree}{$^{***}$}

% ============================================================
\begin{document}

\begin{titlepage}
\begin{center}
\vspace*{1.5cm}

{\LARGE \textbf{IP Policy Communication and Technology Transfer:}}\\[0.4em]
{\LARGE \textbf{Panel Evidence from U.S.\ Research Universities}}

\vspace{1.2cm}

{\large Krishna Sharma\textsuperscript{1},\quad Lily Wang\textsuperscript{2},\quad
Sujan Basnet\textsuperscript{3},\quad Marco Cossco\textsuperscript{4}}

\vspace{0.6cm}
{\normalsize
\textsuperscript{1}Stanford University\quad
\textsuperscript{2}University of California\quad
\textsuperscript{3}MIT\quad
\textsuperscript{4}Carnegie Mellon University}

\vspace{1.2cm}
{\normalsize Working Paper \hspace{2em} September 2026}

\vspace{1.4cm}

\begin{minipage}{0.85\textwidth}
\setstretch{1.3}
\small
\textbf{Abstract.} University intellectual-property policies are governance instruments
that structure how institutions communicate with faculty inventors, technology-transfer
offices, and external partners. Using the Policy Communication Stance Index (\PCSI{})
developed in \citet{sharma2026pcsi}---a sentence-level BERT-based measure constructed
from approximately 87,160 scored sentences drawn from 150 U.S.\ research universities
over 1944--2025---we study whether differences in the communicative character of IP
policies are associated with subsequent technology-transfer activity. Merging \PCSI{}
with AUTM survey outcomes for 253 institutions over 1991--2023, we estimate two-way
fixed-effects panel models with institution and year fixed effects and lagged controls.
The most notable empirical pattern concerns the institution-level log
application-to-disclosure index, defined as the log difference between new patent
applications and invention disclosures: the lag-1 conditional association is
$\hat\beta \approx 1.00$ (SE $\approx 0.50$, $p \approx .047$), and the lag-2
association is $\hat\beta \approx 1.34$ ($p \approx .014$). Associations with raw
patent volume are positive at some longer lags but imprecise at the primary specification.
Broader multiple-testing adjustment removes the nominal significance of the lag-1
result. A falsification exercise against future research expenditure indicates that
common time-varying factors or institutional dynamics correlated with policy revision
cannot be excluded. These patterns are consistent with the view that more accessible
and structured IP policy language corresponds to stronger institutional triage at the
disclosure-to-application margin, though the observational design does not isolate the
mechanism.
\end{minipage}

\vspace{0.8cm}
{\small \textit{Keywords}: technology transfer, intellectual property policy,
text as data, BERT, panel data, university patenting, TTO governance}\\[0.3em]
{\small \textit{JEL}: O31, O34, I23, L31}

\end{center}
\end{titlepage}

% ============================================================
\section{Introduction}
\label{sec:intro}
% ============================================================

The commercialization of academic research depends on a chain of organizational actions
that begins long before a patent application is filed. Faculty members must recognize an
invention as disclosable, decide to approach the technology-transfer office (TTO), and
cooperate with the evaluation and prosecution process. TTOs must triage disclosures
efficiently, allocate resources toward commercially promising inventions, and maintain
productive working relationships with inventor communities that often remain skeptical of
institutional IP ownership. A large literature has documented how the quantity and quality
of TTO resources, the institution's research intensity, royalty-sharing arrangements, and
broader organizational norms shape commercialization outcomes \citep{henderson1998,
mowery2001, thursby2002, thursbythursby2002, siegel2003, siegel2004, dechenaux2011}.
Yet one dimension of the institutional environment has received less systematic attention:
the language through which universities communicate their IP governance to the people who
are expected to operate within it.

University IP policies are formal governance documents. They specify who owns
inventions, what triggers disclosure obligations, how revenue is divided, and what support
inventors can expect. They also do something less easily captured by a binary coding of
policy provisions: they convey, through tone, clarity, and legal register, the institutional
relationship between the university and its faculty inventors. A policy that frames
disclosure as a mandatory compliance obligation communicates something different from one
that presents the TTO as a partner in commercialization, even if the underlying ownership
rules are identical. If that communicative difference shapes how faculty engage with
disclosure and how TTOs organize their evaluation processes, it may be visible in
downstream outcomes at the institution-year level.

This paper studies that question empirically. We build on \citet{sharma2026pcsi}, which
develops and validates the Policy Communication Stance Index (\PCSI{})---a sentence-level
BERT-based measure of communicative stance in university IP policy documents. The
companion paper establishes the measurement and validation; we do not repeat those
details here. Our contribution is to link \PCSI{} to technology-transfer outcomes using
an institution-year panel that merges the policy measure with three decades of AUTM
licensing survey data.

The empirical patterns we document are informative in several respects. First, we
find that the log application-to-disclosure index---the log difference between patent
applications and invention disclosures at the institution-year level---is positively and
statistically associated with \PCSI{} at one- and two-year lags, with the two-year
association ($\hat\beta \approx 1.34$, $p \approx .014$) more precisely estimated than
the one-year result. This outcome is substantively meaningful: it captures whether
institutional triage of disclosed inventions toward patent applications is stronger in
environments with more supportive and structured policy language, abstracting from the
overall scale of research activity. Second, associations with raw patent volume are less
precise at the primary lag, though positive associations emerge at some longer lags for
applications and patents issued. Third, the broader multiple-testing adjustment across
outcomes substantially weakens the nominal significance of the primary result. We discuss
this openly and organize the paper's evidential claims around what the data can and
cannot support.

We treat this as an observational, association-based study throughout. The design
exploits within-institution variation in policy communication over time under the
two-way fixed-effects structure, but we do not have an instrument for \PCSI{} and
do not claim one. A falsification test---regressing future research expenditure on
current \PCSI{}---reveals a positive and marginally significant conditional
association, indicating that reverse causality or correlated time-varying institutional
factors cannot be dismissed. An exploratory event-study diagnostic, described in the
appendix, recovers only three high-confidence policy-revision events with verified
effective dates, and those events exhibit meaningful pre-revision differences in outcomes,
which prevents a credible parallel-trends interpretation. The contribution of this paper
is therefore measurement and pattern-documentation, not causal identification. It
establishes an empirical regularity that motivates a research program rather than
resolving it.

The rest of the paper proceeds as follows. Section~\ref{sec:framework} develops the
conceptual motivation and reviews the relevant literature. Section~\ref{sec:data}
describes the data and measurement. Section~\ref{sec:strategy} presents the empirical
approach. Section~\ref{sec:descriptive} reports descriptive evidence. Section~\ref{sec:results}
presents the panel estimates. Section~\ref{sec:robust} discusses robustness and
interpretation. Section~\ref{sec:discussion} addresses the broader discussion and
limitations. Section~\ref{sec:conclusion} concludes.

% ============================================================
\section{Conceptual Framework and Related Literature}
\label{sec:framework}
% ============================================================

The commercialization of academic inventions has been studied from several angles since
the Bayh-Dole Act of 1980 assigned ownership of federally funded inventions to universities
\citep{mowery2001}. The empirical literature has traced the consequences of that shift through
growth in university patenting \citep{henderson1998, sampat2006}, the determinants of
licensing success \citep{jensenthursby2001, lachschankerman2008}, and the organizational
factors that differentiate high- and low-performing TTOs \citep{thursby2002, siegel2003,
siegel2004, dechenaux2011}. This literature makes clear that resource inputs---TTO staffing,
research funding, institutional prestige---are important predictors of commercialization
output, but also that organizational and relational factors contribute independently.
\citet{owensmithpowell2001} emphasize the role of faculty norms and inter-organizational
networks; \citet{siegel2004} document how the culture of relationships between academic
inventors and TTO professionals influences knowledge transfer in ways that formal
incentive structures do not capture.

Within this organizational strand, the formal policy document has attracted limited
attention as a governance variable. One reason is measurement: IP policies are text, and
converting legal prose into comparable quantitative indicators across institutions and
time periods has not been straightforward. \citet{lachschankerman2008} code royalty-sharing
provisions quantitatively, but that approach captures a single contractual feature rather
than the communicative environment as a whole. The broader study of how formal documents
encode institutional authority and relationship has proceeded primarily in organizational
sociology, where \citet{dimaggio1983} and \citet{meyer1977} established that formal
texts serve legitimacy and signaling functions that may be partially decoupled from
operational behavior. Whether those communicative functions have measurable downstream
consequences for performance is an open empirical question.

Text-as-data methods have expanded the set of institutional features that can be measured
at scale. The approach developed by \citet{loughranmcdonald2011} for financial disclosures
demonstrates that the tone and legal character of formal documents carry information about
institutional behavior beyond the formal content of their rules. \citet{gentzkowkellytaddy2019}
provide a general framework for using text as data in economics, emphasizing that text
scores should be treated as observational proxies subject to the same identification
concerns as any other covariate. \citet{sharma2026pcsi} apply this framework to
university IP policies, constructing \PCSI{} as a sentence-level measure of communicative
stance validated against human ratings and existing policy-quality measures.

The conceptual basis for expecting a relationship between \PCSI{} and
technology-transfer outcomes operates through two plausible channels, neither of which
we are able to test directly. The first is inventor-facing: a policy document that
clearly explains disclosure procedures, defines the scope of covered inventions, describes
what support inventors can expect from the TTO, and presents the institutional relationship
in non-adversarial terms may reduce the perceived transaction cost of disclosure for
faculty members who are otherwise uncertain about what the process entails
\citep{jensenthursby2001, thursbythursby2002}. If policy clarity reduces hesitation at
the disclosure decision, and if the marginal undisclosed invention tends to be one that
the TTO would pursue aggressively, a more communicatively supportive policy could
correspond to a higher ratio of patent applications to disclosures---capturing triage
quality rather than sheer volume.

The second channel operates inside the TTO. IP policies are not only external communications
to faculty; they also constitute formal statements of institutional purpose, scope, and
procedure for TTO staff. A policy that specifies evaluation criteria, triage standards,
and follow-up obligations may coordinate TTO routines in ways that improve the
consistency and efficiency of the disclosure-to-application pipeline. This channel
predicts an association specifically at the application-to-disclosure margin rather than
at the level of raw application counts, because it concerns the institutional conversion
of incoming disclosures rather than the absolute volume of research activity.

We are careful to treat these as plausible interpretive channels rather than established
mechanisms. The data we analyze are aggregate institution-year observations; we do not
observe faculty comprehension, inventor behavior, TTO staffing processes, or the fate
of individual inventions. The findings reported in Section~\ref{sec:results} are
consistent with a TTO-mediated governance interpretation, but they cannot discriminate
between the two channels or rule out alternative explanations. We return to this in
Section~\ref{sec:discussion}.

% ============================================================
\section{Data and Measurement}
\label{sec:data}
% ============================================================

\subsection*{Policy communication measure}

The \PCSI{} measure is developed and validated in \citet{sharma2026pcsi}. Briefly, it
is constructed from sentence-level BERT classifications of IP policy documents collected
from 150 U.S.\ research universities spanning 1944--2025. The sentence classifier is
trained to distinguish communicative stance along dimensions of supportiveness, clarity,
and legal load, and institution-year \PCSI{} scores are obtained by aggregating
sentence-level scores to the policy level and then forward-filling across years until
a revised policy is observed. The measure ranges from approximately 0.3 to 0.7, with
higher values indicating more supportive, clearer, and less legally burdensome language.
The companion paper establishes the reliability and validity of this measure; we take
the institution-year \PCSI{} values as inputs.

One item requires reconciliation between the two papers. \citet{sharma2026pcsi} reports
approximately 87,321 scored sentences in the full corpus. The working corpus used here
contains approximately 87,160 sentences, reflecting the removal of three BYU institution-year
observations (2021--2023) with implausibly long mean sentence lengths consistent with
a parsing artifact, and the correction of one duplicate entry. No institution-level
aggregate \PCSI{} score is materially affected by this adjustment.

\subsection*{Technology-transfer outcomes}

Outcome data come from the AUTM licensing survey, which collects annual institution-level
data on technology-transfer activity from U.S.\ research universities. Our sample covers
253 institutions over 1991--2023, yielding 3,507 institution-year observations. Matching
to the \PCSI{} panel produces 2,564 observations across 150 institutions with at least
one \PCSI{} value. After lagging \PCSI{} and dropping the first year for each institution
in the panel, the baseline estimation sample contains 2,115 observations.

We study seven outcomes from the AUTM survey: new patent applications, total patent
applications, patents issued, invention disclosures, licenses executed, startups formed,
and licensing income. All count outcomes are log-transformed as $\ln(1+y)$ to accommodate
zeros. In addition to these seven outcomes, we construct the \emph{log application-to-disclosure
index}:
\begin{equation}
  A_{it} \;=\; \ln(1 + \text{PatApps}_{it}) \;-\; \ln(1 + \text{Disclosures}_{it}).
  \label{eq:index}
\end{equation}
This index measures the institution-year relationship between patent-application activity
and incoming invention disclosures. It is not a cohort-based conversion rate; disclosures
in year $t$ may generate applications in years $t$ through $t+k$ and the index does not
track individual inventions. What it captures is the institutional tendency, in a given
year, to file patent applications at a relatively high rate compared to the volume of
disclosures arriving at the TTO. A TTO that aggressively pursues the most promising
fraction of its disclosure pipeline will tend to exhibit a higher value of $A_{it}$. We
treat this as our primary outcome because it is the measure most directly responsive to
the triage-efficiency channel described in Section~\ref{sec:framework}, and because it
controls for scale differences in research activity that would otherwise conflate
patent-volume associations with research-intensity differences.

Institutional controls include the log of total research expenditure, the log of
licensing FTEs, and the inventor royalty share, all lagged one period. TLO age is
absorbed by the institution fixed effect and is not separately included.

% ============================================================
\section{Empirical Approach}
\label{sec:strategy}
% ============================================================

We estimate two-way fixed-effects (TWFE) panel regressions of the form:
\begin{equation}
  Y_{it} = \alpha_i + \lambda_t + \beta\,\PCSI{}_{i,t-k} + \mathbf{X}_{i,t-1}'\gamma + \varepsilon_{it},
  \label{eq:twfe}
\end{equation}
where $Y_{it}$ is a technology-transfer outcome for institution $i$ in year $t$,
$\alpha_i$ and $\lambda_t$ are institution and year fixed effects, $\PCSI{}_{i,t-k}$
is the policy communication measure lagged $k$ periods, and $\mathbf{X}_{i,t-1}$ is
the vector of lagged controls. Standard errors are clustered at the institution level
using the \citet{cameron2015} small-sample correction.

The institution fixed effects absorb all time-invariant differences across universities,
including persistent differences in research culture, TTO age, institutional mission, and
faculty composition. The year fixed effects absorb aggregate shocks common to all
universities in a given year, such as changes in patent law, funding cycles, or national
patenting trends. The identifying variation is therefore within-institution
year-to-year deviation in \PCSI{} relative to the institution mean, conditional on
common time trends and the lagged controls.

We use lagged rather than contemporaneous \PCSI{} for two reasons. First, TTO evaluation,
attorney engagement, and patent prosecution involve time delays that make it implausible
for policy communication to be reflected in applications within the same year. Second,
lagging reduces concerns about reverse simultaneity within the same observation period,
though it does not resolve reverse causality operating at longer horizons. We estimate
the primary specification at lag-1 and report results through lag-5 in the appendix
to characterize the full temporal profile. The lag-2 specification is treated as a
secondary result of interest, given the time it typically takes for policy communication
to propagate through the disclosure-to-application pipeline.

A feature of the \PCSI{} panel that affects interpretation is that most policy
observations are carried forward from the most recently available document rather than
freshly observed each year. Approximately 89\% of the matched \PCSI{} observations
reflect a carried-forward value; the mean policy age at observation is 7.8 years
(median 6 years). The within-institution variation in \PCSI{} is therefore
concentrated in the 37 policy-revision events---years in which $|\Delta\PCSI{}|>0.05$,
across 30 institutions---and in the cross-sectional spread of policies for the remaining
institutions. This means the regression is not exploiting fine-grained annual policy
updates but rather comparing trajectories around the comparatively rare moments when
institutions substantially revise their policies, supplemented by persistent
cross-sectional differences.

We estimate equation~(\ref{eq:twfe}) for all eight outcomes. To address multiple
testing, we report Benjamini-Hochberg (BH) adjusted $p$-values alongside raw $p$-values
\citep{bakerlarkerwang2022}. We treat the log application-to-disclosure index as the
primary pre-specified outcome, given the a priori theoretical reasoning in
Section~\ref{sec:framework}; the remaining seven AUTM outcomes are secondary.

% ============================================================
\section{Descriptive Evidence}
\label{sec:descriptive}
% ============================================================

Table~\ref{tab:desc} reports descriptive statistics for the matched sample. The mean
\PCSI{} is 0.40 with a standard deviation of 0.10, and the full range extends from
approximately 0.3 to 0.7. A notable regularity is that 86\% of institution-year
observations fall below the 0.5 neutral midpoint: U.S.\ research university IP policies
are predominantly restrictive in communicative stance. The distribution is rightward
skewed, with a small number of institutions maintaining policies that score in the
upper range. Figure~\ref{fig:motivating} illustrates the cross-sectional density and
the time trend in mean \PCSI{}. Mean communicative stance increased modestly through
the 1990s---consistent with a wave of policy adoption and revision following Bayh-Dole
formalization at many institutions---and has been broadly stable since.

\begin{table}[htbp]
\centering
\caption{Descriptive Statistics}
\label{tab:desc}
\begin{threeparttable}
\begin{tabular}{lrrrrr}
\toprule
Variable & Mean & SD & Min & Max & $N$ \\
\midrule
\multicolumn{6}{l}{\textit{Panel A: Policy Communication Stance Index}} \\[2pt]
\PCSI{} & 0.40 & 0.10 & 0.30 & 0.70 & 2,564 \\
Tone sub-index & 0.00 & 0.50 & $-1.7$ & $4.4$ & 2,564 \\
Clarity sub-index & 0.00 & 0.60 & $-1.9$ & $1.8$ & 2,564 \\
Legal load sub-index & 0.00 & 0.70 & $-1.4$ & $3.2$ & 2,564 \\
\midrule
\multicolumn{6}{l}{\textit{Panel B: Technology Transfer Outcomes}} \\[2pt]
New patent applications & 70.6 & 113.7 & 0 & 1,532 & 3,507 \\
Invention disclosures & 122.4 & 157.7 & 0 & 1,726 & 3,507 \\
Licenses executed & 38.3 & 66.9 & 0 & 1,387 & 3,489 \\
Startups formed & 4.6 & 7.1 & 0 & 102 & 3,245 \\
License income (\$000s) & 11,881 & 46,747 & 0 & 1,258,637 & 3,507 \\
\midrule
\multicolumn{6}{l}{\textit{Panel C: Institutional Characteristics}} \\[2pt]
Total research exp.\ (\$M) & 333.0 & 478.1 & 1.2 & 7,500 & 3,507 \\
Licensing FTEs & 5.5 & 7.3 & 0 & 95 & 3,401 \\
Royalty share (\%) & 39.7 & 8.1 & 15 & 65 & 3,507 \\
TLO age (years) & 23.0 & 15.5 & $-11$ & 97 & 3,507 \\
\bottomrule
\end{tabular}
\begin{tablenotes}
\small
\item \textit{Notes:} Panel A statistics computed on the matched \PCSI{}--AUTM sample.
Panels B and C statistics computed on the full AUTM sample. License income is in
thousands of 2015 dollars.
\end{tablenotes}
\end{threeparttable}
\end{table}

\begin{figure}[htbp]
  \centering
  \includegraphics[width=0.95\textwidth]{figures/fig1_motivating.pdf}
  \caption{\PCSI{} Distribution and Time Trend, 1991--2023}
  \label{fig:motivating}
  \footnotesize
  \textit{Notes:} Left panel: kernel density of institution-year \PCSI{} observations.
  Vertical dashed line marks the 0.5 neutral midpoint; 86\% of observations fall to
  the left. Right panel: annual mean \PCSI{} (solid line) $\pm 1$ SD (shaded band).
\end{figure}

Across the sample, mean annual new patent applications are 71 per institution (SD 114),
while mean annual disclosures are 122 (SD 158). Disclosures therefore substantially
outnumber applications in a typical institution-year, and the log application-to-disclosure
index has a mean near zero with substantial institution-level spread. Institutions with
higher research expenditure and more licensing FTEs tend to be at the upper end of both
the patent-application and disclosure distributions. The raw relationship between \PCSI{}
and the log application-to-disclosure index is positive in the cross section, consistent
with the direction of the within-institution estimates reported below, though the
cross-sectional relationship may conflate \PCSI{} with other institutional characteristics.

Variation in \PCSI{} is predominantly between institutions: institutions differ
substantially and persistently in their communicative stance, while within-institution
changes---concentrated in the 37 policy-revision years identified in the sample---are
comparatively infrequent. This means the TWFE estimator in equation~(\ref{eq:twfe})
weights heavily on cross-sectional differences in the long run and on policy-revision
contrasts in the short run, a feature we return to in the robustness discussion.

% ============================================================
\section{Panel Results}
\label{sec:results}
% ============================================================

Table~\ref{tab:baseline} presents the baseline specification results for log new patent
applications, entering controls sequentially. The \PCSI{} coefficient is positive and
stable across specifications, ranging from 0.69 to 0.62, but not statistically significant
at conventional levels in any column (column (4): $\hat\beta = 0.616$, SE = 0.523,
$p = .241$). Research expenditure is the most important control covariate, entering with
a coefficient of approximately 0.22 and consistent significance, reflecting the expected
positive relationship between institution scale and patent activity.

\begin{table}[htbp]
\centering
\caption{Baseline Specification: Log New Patent Applications}
\label{tab:baseline}
\begin{threeparttable}
\begin{tabular}{lcccc}
\toprule
 & (1) & (2) & (3) & (4) \\
\midrule
\PCSI{}$_{t-1}$ & $0.687$ & $0.581$ & $0.599$ & $0.616$ \\
 & $(0.544)$ & $(0.520)$ & $(0.522)$ & $(0.523)$ \\[4pt]
Lag ln(research exp.) & & $0.225^{**}$ & $0.204^{**}$ & $0.219^{**}$ \\
 & & $(0.098)$ & $(0.103)$ & $(0.108)$ \\[4pt]
Lag ln(licensing FTEs) & & & $0.091$ & $0.091$ \\
 & & & $(0.078)$ & $(0.077)$ \\[4pt]
Lag royalty share & & & & $0.006$ \\
 & & & & $(0.006)$ \\
\midrule
Institution FE & Yes & Yes & Yes & Yes \\
Year FE & Yes & Yes & Yes & Yes \\
$N$ & 2,115 & 2,115 & 2,115 & 2,115 \\
$R^2$ & 0.002 & 0.012 & 0.014 & 0.015 \\
\bottomrule
\end{tabular}
\begin{tablenotes}
\small
\item \textit{Notes:} Dependent variable is $\ln(1+\text{new patent applications})$.
Institution-clustered standard errors in parentheses with \citet{cameron2015} $G-1$
small-sample correction. TLO age is absorbed by institution fixed effects.
$^*p<0.10$, $^{**}p<0.05$, $^{***}p<0.01$.
\end{tablenotes}
\end{threeparttable}
\end{table}

Table~\ref{tab:main} extends the analysis to all eight outcomes and reports multiple-testing
adjustments. Panel~A shows the lag-1 results with the full control set. The most notable
estimate is for the log application-to-disclosure index: $\hat\beta = 1.000$
(SE = 0.500, $p = .047$). The 95\% confidence interval extends from approximately 0.01
to 1.99, confirming positive directionality. For raw patent volumes, the estimates are
positive but imprecise: new applications ($p = .241$), total applications ($p = .144$),
and patents issued ($p = .183$) all fall short of conventional significance thresholds.
Disclosures, licenses, startups, and licensing income are estimated close to zero or
with wide confidence intervals.

Panel~B of Table~\ref{tab:main} reports the Benjamini-Hochberg multiple-testing
adjustment across all eight outcomes at lag-1. The BH-adjusted $p$-value for the log
application-to-disclosure index is 0.374---substantially above conventional thresholds.
The nominal significance of the lag-1 result therefore reflects being the single
lowest $p$-value among eight tested outcomes rather than a finding that withstands
multiple-comparisons adjustment. Panel~C reports lag sensitivity for the two outcomes
with the clearest positive pattern. The lag-2 estimate for the application-to-disclosure
index is $\hat\beta = 1.343$ ($p = .014$), which is both more precise and larger in
magnitude than the lag-1 result. The lag-2 new patent application coefficient is
$\hat\beta = 1.115$ ($p = .042$). At lag-3, the patents-issued coefficient is
$\hat\beta = 1.244$ ($p = .036$). These medium-term associations suggest that the
relationship between \PCSI{} and technology-transfer activity, whatever its source,
operates with a delay consistent with the patent prosecution and evaluation timeline
rather than reflecting contemporaneous institutional dynamics.

The strongest and most consistent pattern in these results concerns the log
application-to-disclosure index. Associations with raw patent volume are less precise
at the primary lag, although positive associations appear at longer lags. No other outcome
exhibits a comparable pattern of repeated significance across lag horizons.

\begin{table}[htbp]
\centering
\caption{Conditional Associations Between \PCSI{} and Technology-Transfer Outcomes}
\label{tab:main}
\begin{threeparttable}
\begin{tabular}{lrrrr}
\toprule
Outcome & $\hat\beta$ & SE & Raw $p$ & $N$ \\
\midrule
\multicolumn{5}{l}{\textit{Panel A: Lag-1 \PCSI{}, full controls}} \\[2pt]
Log application-to-disclosure index & $1.000$ & $0.500$ & $.047$ & 2,115 \\
Ln(new patent applications) & $0.616$ & $0.523$ & $.241$ & 2,115 \\
Ln(total patent applications) & $0.796$ & $0.542$ & $.144$ & 2,097 \\
Ln(patents issued) & $0.763$ & $0.571$ & $.183$ & 2,097 \\
Ln(disclosures) & $-0.384$ & $0.412$ & $.352$ & 2,115 \\
Ln(licenses executed) & $0.570$ & $0.526$ & $.280$ & 2,110 \\
Ln(startups formed) & $0.269$ & $0.533$ & $.614$ & 2,080 \\
Ln(license income) & $0.694$ & $1.126$ & $.538$ & 2,115 \\
\midrule
\multicolumn{5}{l}{\textit{Panel B: Benjamini--Hochberg adjustment (lag-1, all outcomes)}} \\[2pt]
Log application-to-disclosure index & $1.000$ & $0.500$ & $.047$ $[0.374]$ & \\
Ln(new patent applications) & $0.616$ & $0.523$ & $.241$ $[0.447]$ & \\
Ln(total patent applications) & $0.796$ & $0.542$ & $.144$ $[0.447]$ & \\
Ln(patents issued) & $0.763$ & $0.571$ & $.183$ $[0.447]$ & \\
Ln(licenses executed) & $0.570$ & $0.526$ & $.280$ $[0.447]$ & \\
Ln(disclosures) & $-0.384$ & $0.412$ & $.352$ $[0.469]$ & \\
Ln(license income) & $0.694$ & $1.126$ & $.538$ $[0.614]$ & \\
Ln(startups formed) & $0.269$ & $0.533$ & $.614$ $[0.614]$ & \\
\midrule
\multicolumn{5}{l}{\textit{Panel C: Lag sensitivity for selected outcomes}} \\[2pt]
App-to-disclosure, lag-1 & $1.000$ & $0.500$ & $.047$ & 2,115 \\
App-to-disclosure, lag-2 & $1.343$ & $0.540$ & $.014$ & 1,781 \\
App-to-disclosure, lag-3 & $0.918$ & $0.535$ & $.090$ & 1,522 \\
New patent apps, lag-1 & $0.616$ & $0.523$ & $.241$ & 2,115 \\
New patent apps, lag-2 & $1.115$ & $0.545$ & $.042$ & 1,781 \\
Patents issued, lag-3 & $1.244$ & $0.494$ & $.036$ & 1,522 \\
\bottomrule
\end{tabular}
\begin{tablenotes}
\small
\item \textit{Notes:} All specifications include institution and year fixed effects and
lagged controls (log research expenditure, log licensing FTEs, royalty share). Institution-clustered
SEs with \citet{cameron2015} correction. Panel B brackets contain BH-adjusted $p$-values;
the lag-1 result for the primary outcome does not survive multiple-testing adjustment
(BH $p = .374$). $^*p<0.10$, $^{**}p<0.05$, $^{***}p<0.01$.
\end{tablenotes}
\end{threeparttable}
\end{table}

% ============================================================
\section{Robustness and Interpretation}
\label{sec:robust}
% ============================================================

Table~\ref{tab:robust} reports a suite of robustness checks for the two outcomes with
the most interpretable patterns. For the log application-to-disclosure index, the
conditional association survives winsorizing at the 1st and 99th percentiles
($\hat\beta = 1.224$, $p = .038$), restriction to the balanced panel with at least 15
observations ($\hat\beta = 1.046$, $p = .066$), and restriction to Carnegie R1
universities ($\hat\beta = 1.166$, $p = .053$). The lag-2 specification for this
outcome is the single most robust point estimate in the paper ($\hat\beta = 1.343$,
$p = .014$). For log new patent applications, the point estimate is consistently
positive across all robustness checks and the leave-one-out range extends from 0.34
to 0.87 with all 253 institutions yielding positive estimates, though the association
does not reach significance at the primary lag in any specification using conventional
clustered errors. Under HC3 heteroskedasticity-robust errors the coefficient for patent
applications nominally reaches significance ($\hat\beta = 0.616$, SE = 0.308, $p = .046$),
but that reflects a less conservative variance estimator.

Split-sample heterogeneity analyses suggest that associations are stronger at private
universities and non-land-grant institutions, where TTO discretion may be higher. R1
universities show significant associations for both the application-to-disclosure index
and, marginally, for patent applications. These patterns are consistent with a governance
mechanism concentrated in settings where the TTO has a more active role in shaping
disclosure behavior, though the subgroup samples are small and the results should be
treated as exploratory.

The falsification test is the most substantive interpretive challenge. Regressing
future log research expenditure (year $t+1$) on current \PCSI{} yields $\hat\beta = 0.440$
($p = .084$). This positive marginal association is consistent with at least two
alternative explanations. One is reverse causality: universities experiencing expansion
may revise their IP policies in a more supportive direction around the same time that
patent activity increases, generating a spurious association between \PCSI{} and
subsequent outcomes. The other is common-factor confounding: a latent institutional
characteristic---new leadership, a strategic TTO investment, or increased engagement
with industry---may simultaneously drive policy revision, resource growth, and stronger
commercialization performance. Neither explanation can be excluded with the available
data, and both motivate caution about attributing the conditional association to
policy communication per se.

\begin{table}[htbp]
\centering
\caption{Robustness Checks}
\label{tab:robust}
\begin{threeparttable}
\begin{tabular}{lrrrr}
\toprule
Specification & $\hat\beta$ & SE & $p$ & $N$ \\
\midrule
\multicolumn{5}{l}{\textit{Panel A: Log new patent applications}} \\[2pt]
Baseline (lag-1 \PCSI{}) & $0.616$ & $0.523$ & $.241$ & 2,115 \\
Winsorized 1st/99th & $0.841$ & $0.593$ & $.158$ & 2,115 \\
Balanced panel ($\geq 15$ obs) & $0.693$ & $0.595$ & $.246$ & 1,794 \\
R1 universities only & $0.978^*$ & $0.572$ & $.089$ & 1,866 \\
HC3 errors & $0.616^{**}$ & $0.308$ & $.046$ & 2,115 \\
With lagged outcome (Nickell lower bound) & $0.339$ & $0.306$ & $.269$ & 2,115 \\
Falsification: \PCSI{} $\to$ ln(future res.\ exp.) & $0.440^*$ & $0.253$ & $.084$ & 1,773 \\
Leave-one-out range & \multicolumn{4}{l}{$[0.340,\;0.874]$; 100\% positive ($n=253$)} \\
\midrule
\multicolumn{5}{l}{\textit{Panel B: Log application-to-disclosure index}} \\[2pt]
Baseline (lag-1 \PCSI{}) & $1.000^{**}$ & $0.500$ & $.047$ & 2,115 \\
Winsorized 1st/99th & $1.224^{**}$ & $0.586$ & $.038$ & 2,115 \\
Balanced panel ($\geq 15$ obs) & $1.046^*$ & $0.563$ & $.066$ & 1,794 \\
R1 universities only & $1.166^*$ & $0.598$ & $.053$ & 1,866 \\
Lag-2 \PCSI{} & $1.343^{**}$ & $0.540$ & $.014$ & 1,781 \\
\bottomrule
\end{tabular}
\begin{tablenotes}
\small
\item \textit{Notes:} All specifications include institution and year fixed effects and
lagged controls unless otherwise indicated. Falsification test regresses $\ln(\text{research
expenditure}_{t+1})$ on $\PCSI{}_t$ with full controls.
$^*p<0.10$, $^{**}p<0.05$, $^{***}p<0.01$.
\end{tablenotes}
\end{threeparttable}
\end{table}

A circular-shift placebo test further characterizes the null distribution of the patent-application
coefficient. Rotating each institution's \PCSI{} time series by a random number of periods
and re-estimating the baseline regression 2,000 times yields a placebo $p$-value of 0.193
for new patent applications: the observed coefficient is not unusual in the permutation
distribution, consistent with the absence of conventional significance in that outcome.

The sub-index decomposition reported in Appendix~\ref{app:channel} finds no single
dimension of policy communication that drives the primary association. Tone, clarity,
and legal load each enter with imprecise and inconsistent estimates when substituted for
the composite \PCSI{} or included jointly. This could reflect collinearity among the
sub-dimensions, noise in the individual sub-indices, or genuinely diffuse source of
the composite association.

% ============================================================
\section{Discussion and Limitations}
\label{sec:discussion}
% ============================================================

The results presented above establish a positive conditional association between
institution-level \PCSI{} and the log application-to-disclosure index at one- and
two-year lags, with the two-year estimate being both stronger and more precisely
identified. This pattern is consistent with the view that universities whose IP policies
communicate more accessibly and supportively exhibit stronger triage efficiency at the
disclosure-to-application margin. The evidence is, however, observational, and several
important limitations bear on the interpretation.

The identification challenge is serious. The institution fixed effects remove
time-invariant differences across universities, but they do not account for
time-varying institutional dynamics that could simultaneously produce policy revisions
and performance changes. The falsification test reveals that \PCSI{} is marginally
associated with future research expenditure, which is consistent with reverse causality
or with a common driver such as new administrative leadership or a strategic TTO
investment. Universities experiencing resource expansions may also revise their IP
policies in a more welcoming direction, creating a spurious association between the
measure and subsequent commercialization outcomes. We cannot identify which explanation
is operative, and the available data do not include the fine-grained administrative or
personnel records that would be needed to test this more carefully.

The policy-revision structure of the data also limits inference. Approximately 89\%
of \PCSI{} observations in the estimation sample reflect carried-forward values from
an institution's most recently available policy, with a mean policy age of 7.8 years
at observation. The within-institution variation that enters the fixed-effects estimator
is therefore predominantly concentrated in the 37 institution-years of substantial
policy revision, and at longer policy ages the measure may reflect a governance
environment that has evolved away from the formal document. Whether a decade-old policy
document meaningfully describes current TTO practice is uncertain.

Measurement error in \PCSI{} enters through the sentence classification process
itself, through aggregation from sentence to institution-year, and through the
forward-fill assumption. Attenuation toward zero is the expected direction of the
resulting bias. The fact that the coefficients are positive and of moderate magnitude
in most specifications is therefore consistent with meaningful signal in the measure,
even if the true relationship is larger.

The AUTM survey introduces its own selection concern. Participation is voluntary,
and institutions with more active TTOs are more likely to respond. If the positive
association between \PCSI{} and the application-to-disclosure index is stronger among
active institutions, our estimates may not generalize to the full population of U.S.\
research universities. Additionally, the log application-to-disclosure index is an
accounting identity at the institution-year level, not a record of individual invention
trajectories. Changes in the index can reflect shifts in disclosure norms, changes in
patent-filing thresholds, variation in TTO capacity, or strategic responses to changes
in patent law, in addition to any effect of policy communication. The institutional
aggregate masks substantial heterogeneity in inventor- and invention-level processes
that we are unable to observe.

The event-study diagnostic described in Appendix~\ref{app:eventdiag} illustrates
the difficulty of treating policy revisions as quasi-experimental shocks. A
high-confidence audit recovers only three policy-revision events with verified
effective dates and AUTM survey continuity---from Louisiana State University, Oklahoma
State University, and the University of New Orleans. All three display substantial
pre-revision differences in outcomes relative to comparison institutions, undermining
the parallel-trends assumption that event-study interpretation requires \citep{callaway2021,
roth2022, rambachanroth2023}. We include the diagnostic to document the structure of
available variation rather than to supply a causal estimate, and we do not treat it
as corroborating the panel results.

For policy, the findings are suggestive rather than prescriptive. The cross-sectional
regularity that 86\% of U.S.\ research university IP policies score below the
communicative-stance neutral midpoint indicates that most institutions maintain policies
that are predominantly legalistic and restrictive in tone. Whether this reflects deliberate
risk management, legal drafting conventions, or institutional inertia is unclear. The
positive association between more supportive policy communication and the application-to-disclosure
margin---however cautiously it must be qualified---is consistent with the view that the
communicative character of IP policies is not institutionally inconsequential. Universities
reviewing their IP governance documents may find it worthwhile to consider whether the
language they use clearly communicates disclosure procedures, explains TTO evaluation
timelines, and presents the institution as a supportive partner in commercialization.
Establishing whether such revisions actually change faculty behavior or TTO performance
requires designs that go beyond what is available here.

% ============================================================
\section{Conclusion}
\label{sec:conclusion}
% ============================================================

This paper documents a positive association between the communicative character of
university IP policies, measured by the text-based \PCSI{} from \citet{sharma2026pcsi},
and the log application-to-disclosure index at one- and two-year lags in an institution-year
panel. The most robust empirical pattern is at the two-year horizon ($\hat\beta \approx 1.34$,
$p \approx .014$), consistent with the time required for policy communication to propagate
through the TTO evaluation and patent prosecution process. The association does not
survive Benjamini-Hochberg correction at lag-1, and the falsification test indicates
that common time-varying factors cannot be excluded. The paper's contribution is
therefore to establish this empirical regularity---and to document the conditions under
which it holds and fails---rather than to resolve the causal question.

The evidence motivates, but does not settle, whether communication itself changes
commercialization outcomes. Future research that can leverage verified policy adoption
dates with sharp institutional comparisons, TTO leadership transitions, archived effective
dates, invention-level records, or faculty survey data would be considerably better
placed to address the identification challenge. The measurement infrastructure developed
in \citet{sharma2026pcsi} provides the outcome variable for such designs; the remaining
requirement is the source of identifying variation.

% ============================================================
% References
% ============================================================

\bibliographystyle{aer}
\begin{thebibliography}{99}

\bibitem[Sharma et~al.(2026)]{sharma2026pcsi}
Sharma, K., Wang, L., Basnet, S., \& Cossco, M. (2026). Communication as governance:
Measuring communicative stance in university intellectual property policies.
\textit{Working Paper}, September 2026.

\bibitem[AUTM(2023)]{autm2023}
AUTM. (2023). \textit{AUTM Licensing Activity Survey: FY2022}. Association of
University Technology Managers.

\bibitem[Baker et~al.(2022)]{bakerlarkerwang2022}
Baker, A.~C., Larcker, D.~F., \& Wang, C.~C.~Y. (2022). How much should we trust
staggered difference-in-differences estimates? \textit{Journal of Financial Economics},
144(2), 370--395.

\bibitem[Boru\v{s}yak et~al.(2024)]{borusyakjaravel2024}
Boru\v{s}yak, K., Jaravel, X., \& Spiess, J. (2024). Revisiting event study designs:
Robust and efficient estimation. \textit{Review of Economic Studies}, 91(6), 3253--3285.

\bibitem[Callaway and Sant'Anna(2021)]{callaway2021}
Callaway, B., \& Sant'Anna, P.~H.~C. (2021). Difference-in-differences with multiple
time periods. \textit{Journal of Econometrics}, 225(2), 200--230.

\bibitem[Cameron et~al.(2015)]{cameron2015}
Cameron, A.~C., \& Miller, D.~L. (2015). A practitioner's guide to cluster-robust
inference. \textit{Journal of Human Resources}, 50(2), 317--372.

\bibitem[de~Chaisemartin and D'Haultf{\oe}uille(2020)]{dechaisemartin2020}
de~Chaisemartin, C., \& D'Haultf{\oe}uille, X. (2020). Two-way fixed effects estimators
with heterogeneous treatment effects. \textit{American Economic Review}, 110(9),
2964--2996.

\bibitem[Dechenaux et~al.(2011)]{dechenaux2011}
Dechenaux, E., Goldfarb, B., Shane, S., \& Thursby, M. (2011). Appropriability and
the adoption of new technologies: The case of university licensing. \textit{Management
Science}, 57(7), 1205--1220.

\bibitem[Devlin et~al.(2019)]{devlin2019}
Devlin, J., Chang, M.-W., Lee, K., \& Toutanova, K. (2019). BERT: Pre-training of
deep bidirectional transformers for language understanding. In \textit{Proceedings of
NAACL-HLT 2019}, 4171--4186.

\bibitem[DiGregorio and Shane(2003)]{digregorio2003}
DiGregorio, D., \& Shane, S. (2003). Why do some universities generate more start-ups
than others? \textit{Research Policy}, 32(2), 209--227.

\bibitem[DiMaggio and Powell(1983)]{dimaggio1983}
DiMaggio, P.~J., \& Powell, W.~W. (1983). The iron cage revisited: Institutional
isomorphism and collective rationality in organizational fields. \textit{American
Sociological Review}, 48(2), 147--160.

\bibitem[Farre-Mensa et~al.(2020)]{farremensa2020}
Farre-Mensa, J., Hegde, D., \& Ljungqvist, A. (2020). What is a patent worth? Evidence
from the U.S.\ patent ``lottery.'' \textit{Journal of Finance}, 75(2), 639--682.

\bibitem[Goodman-Bacon(2021)]{goodmanbacon2021}
Goodman-Bacon, A. (2021). Difference-in-differences with variation in treatment timing.
\textit{Journal of Econometrics}, 225(2), 254--277.

\bibitem[Henderson et~al.(1998)]{henderson1998}
Henderson, R., Jaffe, A.~B., \& Trajtenberg, M. (1998). Universities as a source of
commercial technology: A detailed analysis of university patenting, 1965--1988.
\textit{Review of Economics and Statistics}, 80(1), 119--127.

\bibitem[Jaffe(1993)]{jaffe1993}
Jaffe, A.~B. (1993). The US patent system in transition: Policy innovation and the
innovation process. \textit{Research Policy}, 29(4--5), 531--557.

\bibitem[Jensen and Thursby(2001)]{jensenthursby2001}
Jensen, R., \& Thursby, M. (2001). Proofs and prototypes for sale: The licensing of
university inventions. \textit{American Economic Review}, 91(1), 240--259.

\bibitem[Lach and Schankerman(2008)]{lachschankerman2008}
Lach, S., \& Schankerman, M. (2008). Incentives and invention in universities.
\textit{RAND Journal of Economics}, 39(2), 403--433.

\bibitem[Lerner and Seru(2022)]{lernerseru2022}
Lerner, J., \& Seru, A. (2022). The use and misuse of patent data: Issues for finance
and beyond. \textit{Review of Financial Studies}, 35(6), 2667--2704.

\bibitem[Loughran and McDonald(2011)]{loughranmcdonald2011}
Loughran, T., \& McDonald, B. (2011). When is a liability not a liability? Textual
analysis, dictionaries, and 10-Ks. \textit{Journal of Finance}, 66(1), 35--65.

\bibitem[Gentzow et~al.(2019)]{gentzkowkellytaddy2019}
Gentzkow, M., Kelly, B.~T., \& Taddy, M. (2019). Measuring group differences in
high-dimensional choices: Method and application to congressional speech.
\textit{Econometrica}, 87(4), 1307--1340.

\bibitem[Gentzkow et~al.(2019)]{gentzkow2019}
Gentzkow, M., Kelly, B.~T., \& Taddy, M. (2019). Text as data. \textit{Journal of
Economic Literature}, 57(3), 535--574.

\bibitem[Meyer and Rowan(1977)]{meyer1977}
Meyer, J.~W., \& Rowan, B. (1977). Institutionalized organizations: Formal structure
as myth and ceremony. \textit{American Journal of Sociology}, 83(2), 340--363.

\bibitem[Mowery et~al.(2001)]{mowery2001}
Mowery, D.~C., Nelson, R.~R., Sampat, B.~N., \& Ziedonis, A.~A. (2001).
The growth of patenting and licensing by U.S.\ universities: An assessment of the
effects of the Bayh-Dole Act of 1980. \textit{Research Policy}, 30(1), 99--119.

\bibitem[Owen-Smith and Powell(2001)]{owensmithpowell2001}
Owen-Smith, J., \& Powell, W.~W. (2001). To patent or not: Faculty decisions and
institutional success at technology transfer. \textit{Journal of Technology Transfer},
26(1), 99--114.

\bibitem[Rambachan and Roth(2023)]{rambachanroth2023}
Rambachan, A., \& Roth, J. (2023). A more credible approach to parallel trends.
\textit{Review of Economic Studies}, 90(5), 2555--2591.

\bibitem[Roth(2022)]{roth2022}
Roth, J. (2022). Pre-test with caution: Event-study estimates after testing for parallel
trends. \textit{American Economic Review: Insights}, 4(3), 305--322.

\bibitem[Sampat(2006)]{sampat2006}
Sampat, B.~N. (2006). Patenting and US academic research in the 20th century: The world
before and after Bayh-Dole. \textit{Research Policy}, 35(6), 772--789.

\bibitem[Siegel et~al.(2003)]{siegel2003}
Siegel, D.~S., Waldman, D., Atwater, L., \& Link, A.~N. (2003). Commercial knowledge
transfers from universities to firms: Improving the effectiveness of university-industry
collaboration. \textit{Journal of High Technology Management Research}, 14(1), 111--133.

\bibitem[Siegel et~al.(2004)]{siegel2004}
Siegel, D.~S., Waldman, D.~A., Atwater, L.~E., \& Link, A.~N. (2004). Toward a model
of the effective transfer of scientific knowledge from academicians to practitioners.
\textit{Journal of Engineering and Technology Management}, 21(1), 115--142.

\bibitem[Suchman(1995)]{suchman1995}
Suchman, M.~C. (1995). Managing legitimacy: Strategic and institutional approaches.
\textit{Academy of Management Review}, 20(3), 571--610.

\bibitem[Sun and Abraham(2021)]{sunabr2021}
Sun, L., \& Abraham, S. (2021). Estimating dynamic treatment effects in event studies
with heterogeneous treatment effects. \textit{Journal of Econometrics}, 225(2), 175--199.

\bibitem[Thursby and Kemp(2002)]{thursby2002}
Thursby, J.~G., \& Kemp, S. (2002). Growth and productive efficiency of university
intellectual property licensing. \textit{Research Policy}, 31(1), 109--124.

\bibitem[Thursby and Thursby(2002)]{thursbythursby2002}
Thursby, J.~G., \& Thursby, M.~C. (2002). Who is selling the ivory tower? Sources of
growth in university licensing. \textit{Management Science}, 48(1), 90--104.

\end{thebibliography}

% ============================================================
\appendix

\section{Lag Sensitivity: Full Results}
\label{app:lagsens}

\begin{table}[htbp]
\centering
\caption{Lag Sensitivity: All Outcomes, Contemporaneous through Lag-5}
\label{tab:lagsens}
\begin{threeparttable}
\begin{tabular}{lrrrrrr}
\toprule
Outcome & Contemp. & Lag 1 & Lag 2 & Lag 3 & Lag 4 & Lag 5 \\
\midrule
Log app-to-disclosure & $0.720$ & $1.000^{**}$ & $1.343^{**}$ & $0.918^*$ & $0.411$ & $0.382$ \\
Ln(new patent apps) & $0.436$ & $0.616$ & $1.115^{**}$ & $0.584$ & $0.195$ & $0.443$ \\
Ln(total patent apps) & $0.732$ & $0.796$ & $0.905$ & $0.498$ & $0.411$ & $0.683$ \\
Ln(patents issued) & $0.496$ & $0.763$ & $1.114^*$ & $1.244^{**}$ & $1.446^{**}$ & $1.753^{***}$ \\
Ln(disclosures) & $-0.284$ & $-0.384$ & $-0.229$ & $-0.334$ & $-0.216$ & $0.061$ \\
Ln(licenses executed) & $0.596$ & $0.570$ & $0.214$ & $0.102$ & $0.342$ & $-0.743$ \\
Ln(startups) & $0.572$ & $0.269$ & $-0.259$ & $-0.619$ & $-0.203$ & $-0.152$ \\
Ln(license income) & $0.654$ & $0.694$ & $0.989$ & $0.215$ & $0.382$ & $0.783$ \\
\bottomrule
\end{tabular}
\begin{tablenotes}
\small
\item \textit{Notes:} Each cell is the \PCSI{} coefficient from a separate TWFE regression
at the indicated lag with full controls and institution-clustered SEs.
$^*p<0.10$, $^{**}p<0.05$, $^{***}p<0.01$.
\end{tablenotes}
\end{threeparttable}
\end{table}

\section{Sub-Index Correlation Matrix and Decomposition}
\label{app:channel}

\begin{table}[htbp]
\centering
\caption{Sub-Index Correlation Matrix (Lagged One Year)}
\label{tab:corrmatrix}
\begin{tabular}{lrrrr}
\toprule
 & \PCSI{} & Tone & Clarity & Legal Load \\
\midrule
\PCSI{} & 1.000 & & & \\
Tone & 0.432 & 1.000 & & \\
Clarity & 0.037 & 0.153 & 1.000 & \\
Legal Load & 0.190 & $-0.026$ & $-0.166$ & 1.000 \\
\bottomrule
\end{tabular}
\begin{tablenotes}
\small
\item \textit{Notes:} Correlations computed on the estimation sample ($N=2{,}115$).
Clarity is nearly orthogonal to the composite \PCSI{}.
\end{tablenotes}
\end{table}

\begin{table}[htbp]
\centering
\caption{Sub-Index Decomposition}
\label{tab:subindex}
\begin{threeparttable}
\begin{tabular}{lrrrr}
\toprule
Sub-index & $\hat\beta$ & SE & $p$ & $N$ \\
\midrule
\multicolumn{5}{l}{\textit{Panel A: Sequential regressions, log new patent applications}} \\[2pt]
Tone index & $0.050$ & $0.077$ & $.513$ & 2,115 \\
Clarity index & $-0.141$ & $0.092$ & $.124$ & 2,115 \\
Legal load index & $0.008$ & $0.070$ & $.906$ & 2,115 \\
\midrule
\multicolumn{5}{l}{\textit{Panel B: Sequential regressions, log application-to-disclosure index}} \\[2pt]
Tone index & $0.009$ & $0.062$ & $.887$ & 2,115 \\
Clarity index & $-0.082$ & $0.070$ & $.241$ & 2,115 \\
Legal load index & $0.029$ & $0.062$ & $.649$ & 2,115 \\
\bottomrule
\end{tabular}
\begin{tablenotes}
\small
\item \textit{Notes:} Each row is a separate regression substituting the indicated
sub-index for \PCSI{} with full controls. No sub-index reaches significance for
either outcome; the source of the composite association is not clearly localized
in a single dimension.
\end{tablenotes}
\end{threeparttable}
\end{table}

\section{Exploratory Policy-Revision Diagnostic}
\label{app:eventdiag}

This appendix documents an exploratory analysis of outcome trajectories around policy
revisions identified in the \PCSI{} panel. It is not an event study in the causal
sense and should not be interpreted as one.

The full panel contains 37 institution-years in which $|\Delta\PCSI{}|>0.05$, across
30 institutions. A high-confidence audit that requires an independently verified
effective date from a public document (board minutes, press release, or archived policy
repository) and AUTM survey continuity in the surrounding years recovers three qualifying
events: Louisiana State University, Oklahoma State University, and the University of
New Orleans. All three exhibit substantial pre-revision differences in outcomes relative
to comparison institutions. Under any of the leading difference-in-differences estimators
for staggered adoption \citep{callaway2021, sunabr2021, dechaisemartin2020,
goodmanbacon2021}, a design in which the parallel-trends assumption is violated in this
manner does not support a causal interpretation of the post-revision coefficients
\citep{roth2022, rambachanroth2023}.

Figure~\ref{fig:pretrend} shows the event-study coefficients for new patent applications
relative to the revision year. The positive pre-revision trend visible in the figure
is the diagnostic evidence that the identified revising institutions were already on a
divergent trajectory, which prevents attribution of the post-revision pattern to the
policy revision itself. We include this figure as a transparency exercise, not as
evidence that the panel estimates in Section~\ref{sec:results} are causally identified.

\begin{figure}[htbp]
  \centering
  \includegraphics[width=0.85\textwidth]{figures/fig2_pretrend.pdf}
  \caption{Exploratory Event-Study Diagnostic: Log New Patent Applications}
  \label{fig:pretrend}
  \footnotesize
  \textit{Notes:} Event-study coefficients and 95\% confidence intervals normalized
  to zero at event-time $-1$. The positive pre-revision trend indicates that parallel
  trends is not credible in this sample. These estimates are descriptive and do not
  support a causal interpretation.
\end{figure}

\end{document}
